\section{Introduction}

Language model agents have demonstrated remarkable proficiency at executing isolated, short-horizon tasks.
In software engineering, agents can resolve individual GitHub issues by localizing bugs and producing patches~\citep{jimenez2024swebench}.
In customer service, they can handle user requests by consulting policy documents and invoking APIs~\citep{yao2024taubench}.
In web navigation, they can complete multi-step browsing tasks across realistic websites~\citep{zhou2024webarena}.
Yet these successes share a common structure: the agent receives a well-defined objective, acts over a brief horizon, and receives relatively immediate feedback on whether it succeeded.
%
What we recognize as intelligence in practice---what distinguishes a capable strategist from a competent executor---often demands something qualitatively different: sustained reasoning under uncertainty over long time horizons, in environments that shift beneath the agent's feet.
%
Such settings are partially observable, requiring the agent to infer latent state from noisy, incomplete signals; non-stationary, as the dynamics themselves evolve over time; and stochastic, so that identical actions yield different outcomes across episodes.
%
Crucially, these environments present the agent with a large, heterogeneous action space---spanning analysis, resource allocation, long-term investment, negotiation, and operational control---and demand that these diverse capabilities be orchestrated into a coherent strategy rather than exercised in isolation.
%
Operating a business is a canonical instance.
A CEO must interpret ambiguous market signals to infer customer sentiment that is never directly revealed, adapt to shifting preferences and macroeconomic conditions that alter the rules of the game mid-play, and simultaneously coordinate pricing, capacity planning, marketing allocation, R\&D investment, and enterprise negotiations---all while absorbing the consequences of decisions made weeks or months earlier.

Existing agent benchmarks do not evaluate this type of capability.
SWE-Bench~\citep{jimenez2024swebench}, WebArena~\citep{zhou2024webarena}, AgentBench~\citep{liu2024agentbench}, and GAIA~\citep{mialon2024gaia} present agents with episodic tasks in largely stationary environments, where the horizon spans minutes to hours and the outcome of each action is promptly revealed.
Even benchmarks that involve richer interaction---such as $\tau$-bench~\citep{yao2024taubench}, which formulates customer service as a POMDP with stochastic user behavior---are scoped to single conversations lasting a handful of turns, with no persistent consequences across episodes.
More broadly, these benchmarks can be characterized as short-horizon, near-fully-observable problems with small action spaces and dense reward signals: settings where reactive execution suffices and long-term strategic planning is unnecessary.
%
VendingBench~\citep{backlund2025vendingbench} takes an important step toward long-horizon evaluation by simulating a multi-day vending machine operation.
However, the environment is deliberately simplified: the agent manages a single four-slot machine with no competitors, near-deterministic demand, and a narrow action space limited to restocking and pricing two to three products.
The strategic depth is minimal---the optimal policy is largely apparent---and the environment lacks the non-stationarity, partial observability over latent market dynamics, and multi-dimensional tradeoffs that characterize genuinely complex decision-making.

We introduce \textsc{CEOBench}, a benchmark that evaluates language model agents as strategic decision-makers in a richly simulated, long-horizon environment.
In \textsc{CEOBench}, an agent operates a subscription-based software company over 365 simulated days, making daily decisions across six interdependent domains: pricing and packaging, infrastructure capacity, marketing and customer acquisition, research and development, enterprise sales negotiations, and financial management.
The agent has access to over 30 tools and a relational database with 34 tables, but critical state variables---individual customer satisfaction, group reputation, churn probabilities, and the true parameters governing preference drift---remain hidden.
The environment is non-stationary: customer preferences drift monthly, macroeconomic conditions follow a stochastic mean-reverting process published with delay, and competitor events trigger exogenous shocks to expected quality.
Feedback is sparse and delayed: revenue arrives on 30-day billing cycles, R\&D projects take 5--30 days to complete, and the reputational consequences of a pricing mistake may not manifest for weeks.
The agent's score---final cash balance at the end of the episode---reflects the cumulative impact of hundreds of interdependent decisions, with bankruptcy (cash below zero) ending the episode early.

\textsc{CEOBench} is designed to stress-test capabilities that existing benchmarks leave largely unexamined.
%
\emph{Long-horizon credit assignment.}
With 365 daily decision points and 30-day billing cycles, the agent must connect current outcomes to actions taken weeks earlier---a challenge that grows combinatorially with the horizon and resists the trial-and-error strategies that succeed in short-episode settings.
%
\emph{Inference under partial observability.}
The agent never observes customer satisfaction, reputation scores, or preference parameters directly; it must construct a mental model of latent market state from aggregate metrics, financial trends, and social media sentiment, then act on beliefs rather than ground truth.
%
\emph{Adaptation to non-stationarity.}
Monthly preference drift, macroeconomic cycles, and stochastic competitor events mean that a strategy optimal in month one may fail by month six.
The agent must detect distributional shift from indirect evidence and adjust accordingly, without ever being told that the environment has changed.
%
\emph{Multi-capability coordination.}
Unlike benchmarks where the agent exercises a single skill (coding, browsing, conversation), \textsc{CEOBench} requires simultaneous orchestration of pricing strategy, capacity planning, marketing allocation, R\&D portfolio management, and multi-turn enterprise negotiations---with tight coupling between domains, so that a capacity shortfall degrades service quality, which erodes reputation, which reduces acquisition, which lowers revenue, which constrains future investment.
%
Together, these properties make \textsc{CEOBench} a challenging testbed for evaluating whether language model agents can move beyond competent task execution toward the kind of sustained, adaptive, strategic reasoning that complex real-world environments demand.
